\scanpage{12}
\mspage{- 3 -}\hfill n\handwritten{°354}

\uline{L'anneau} $A$, \uline{muni de cette graduation}, \uline{est alors appelé anneau gradué de type}
$\handwritten{\Delta}$ .

PROPOSITION 1.- \uline{Si tout élément de} $\handwritten{\Delta}$ \uline{est régulier}, \uline{et si} $(A_{\handwritten{\lambda}})$ \uline{est une graduation de type}
$\handwritten{\Delta}$ \uline{compatible avec la structure d'un anneau} $A$, $A_{\unsure{o|0}}$ \uline{est un}
\uline{sous-anneau de} $A$ (\uline{et en particulier on a} $1\handwritten{\in} A_o$)

Comme $A_o A_o\handwritten{\subset} A_o$ par définition, il suffit de prouver la seconde assertion.
Soit $1=\handwritten{\sum}\limits_{\handwritten{\lambda\in\Delta}}e_{\handwritten{\lambda}}$ la décomposition de 1 en ses composantes homogènes. Si
$x\handwritten{\in}A_\mu$, on a $x=x.1=\handwritten{\sum}\limits_{\handwritten{\lambda\in\Delta}}xe_{\handwritten{\lambda}}$ ; comparant les composantes de degré $\mu$,
on obtient (puisque la relation $\mu+\handwritten{\lambda}=\mu$ entraîne $\handwritten{\lambda}=0$) $x=xe_o$. Cette
relation étant vraie pour tout élément homogène de $A$, est vraie pour tout
$x\handwritten{\in}A$ ; en particulier $1=1.e_o=e_o\handwritten{\in} A_o$ .

DEFINITION 3.- \uline{Soient} $A$ \uline{un anneau gradué de type} $\handwritten{\Delta}$ , $M$ \uline{un} $A$-\uline{module à gauche}
(resp. \uline{à droite}) ; \uline{on dit qu'une graduation} $(M_{\handwritten{\lambda}})$ \uline{de type} $\handwritten{\Delta}$ \uline{est compatible}
\uline{avec la structure de} $A$-\uline{module de} $M$ \uline{si l'on a}
\begingroup
\makeatletter\tagsleft@true\makeatother
\[
 A_{\handwritten{\lambda}} M_\mu \handwritten{\subset} M_{\handwritten{\lambda}+\mu}
 \qquad (\text{resp. }M_\mu A_{\handwritten{\lambda}}\handwritten{\subset} M_{\handwritten{\lambda}+\mu})
 \tag*{(2)}
\]
\endgroup
\uline{quels que soient} $\handwritten{\lambda}$, $\mu$ \uline{dans} $\handwritten{\Delta}$ . \uline{Le module} $M$, \uline{muni de cette graduation}, \uline{est}
\uline{alors appelé module gradué à gauche} (resp. \uline{à droite}) \uline{de type} $\handwritten{\Delta}$ .

Lorsque les éléments de $\handwritten{\Delta}$ sont réguliers, il résulte de (2) et de la
prop.1 que les $M_{\handwritten{\lambda}}$ sont des $A_o$-\uline{modules}.

Il est clair que si $A$ est un anneau gradué de type $\handwritten{\Delta}$ , le $A$-module à gauche $A_s$ (resp. le $A$-module à droite $A_d$) est gradué de type $\handwritten{\Delta}$ .

\uline{Exemples}.- 4) Sur un anneau quelconque $A$, la graduation triviale de type
est compatible avec la structure d'anneau. Si $A$ est gradué par la graduation
triviale, pour qu'une graduation $(M_{\handwritten{\lambda}})$ de type $\handwritten{\Delta}$ sur un $A$-module $M$ soit compatible avec la structure de $A$-module, il faut et il suffit que les $M_{\handwritten{\lambda}}$
soient des \uline{sous-modules} de $M$.
% archive sheet number: 11
